\documentclass[10pt,a4paper]{amsart}
\usepackage{amsmath,amssymb,amsthm,graphicx}
\usepackage{verbatim}
\newenvironment{reference}{\comment}{\endcomment}
%\newenvironment{reference}{}{}
\newenvironment{slogan}{\comment}{\endcomment}
\newenvironment{history}{\comment}{\endcomment}

% For commutative diagrams we use Xy-pic
\usepackage[all]{xy}

% We use 2cell for 2-commutative diagrams.
\xyoption{2cell}
% Theorem environments.
%
\theoremstyle{plain}
\newtheorem{theorem}[subsection]{Theorem}
\newtheorem{proposition}[subsection]{Proposition}
\newtheorem{lemma}[subsection]{Lemma}

\theoremstyle{definition}
\newtheorem{definition}[subsection]{Definition}
\newtheorem{example}[subsection]{Example}
\newtheorem{exercise}[subsection]{Exercise}
\newtheorem{situation}[subsection]{Situation}

\theoremstyle{remark}
\newtheorem{remark}[subsection]{Remark}
\newtheorem{remarks}[subsection]{Remarks}

\numberwithin{equation}{subsection}

% Macros
%
\def\lim{\mathop{\mathrm{lim}}\nolimits}
\def\colim{\mathop{\mathrm{colim}}\nolimits}
\def\Spec{\mathop{\mathrm{Spec}}}
\def\Hom{\mathop{\mathrm{Hom}}\nolimits}
\def\Ext{\mathop{\mathrm{Ext}}\nolimits}
\def\SheafHom{\mathop{\mathcal{H}\!\mathit{om}}\nolimits}
\def\SheafExt{\mathop{\mathcal{E}\!\mathit{xt}}\nolimits}
\def\Sch{\mathit{Sch}}
\def\Mor{\mathop{\mathrm{Mor}}\nolimits}
\def\Ob{\mathop{\mathrm{Ob}}\nolimits}
\def\Sh{\mathop{\mathit{Sh}}\nolimits}
\def\NL{\mathop{N\!L}\nolimits}
\def\CH{\mathop{\mathrm{CH}}\nolimits}
\def\proetale{{pro\text{-}\acute{e}tale}}
\def\etale{{\acute{e}tale}}
\def\QCoh{\mathit{QCoh}}
\def\Ker{\mathop{\mathrm{Ker}}}
\def\Im{\mathop{\mathrm{Im}}}
\def\Coker{\mathop{\mathrm{Coker}}}
\def\Coim{\mathop{\mathrm{Coim}}}

% Boxtimes
%
\DeclareMathSymbol{\boxtimes}{\mathbin}{AMSa}{"02}

%
% Macros for moduli stacks/spaces
%
\def\QCohstack{\mathcal{QC}\!\mathit{oh}}
\def\Cohstack{\mathcal{C}\!\mathit{oh}}
\def\Spacesstack{\mathcal{S}\!\mathit{paces}}
\def\Quotfunctor{\mathrm{Quot}}
\def\Hilbfunctor{\mathrm{Hilb}}
\def\Curvesstack{\mathcal{C}\!\mathit{urves}}
\def\Polarizedstack{\mathcal{P}\!\mathit{olarized}}
\def\Complexesstack{\mathcal{C}\!\mathit{omplexes}}
% \Pic is the operator that assigns to X its picard group, usage \Pic(X)
% \Picardstack_{X/B} denotes the Picard stack of X over B
% \Picardfunctor_{X/B} denotes the Picard functor of X over B
\def\Pic{\mathop{\mathrm{Pic}}\nolimits}
\def\Picardstack{\mathcal{P}\!\mathit{ic}}
\def\Picardfunctor{\mathrm{Pic}}
\def\Deformationcategory{\mathcal{D}\!\mathit{ef}}
\usepackage{cleveref}
\crefname{originallemma}{original lemma}{original lemmas}
\Crefname{originallemma}{Original lemma}{Original lemmas}
\makeatletter
\newlabel{categories-lemma-exact-adjoint}{{0039}{source}{}{sourcepointer-1}{}}
\newlabel{categories-lemma-exact-adjoint@cref}{{[originallemma][0][]0039}{[1][0][]source}{}{}{}}
\makeatother
\crefname{originallemma}{original lemma}{original lemmas}
\Crefname{originallemma}{Original lemma}{Original lemmas}
\makeatletter
\newlabel{lemma-exact-functor}{{010N}{source}{}{sourcepointer-2}{}}
\newlabel{lemma-exact-functor@cref}{{[originallemma][0][]010N}{[1][0][]source}{}{}{}}
\makeatother
\title{A functor with an exact left adjoint preserves injectives}
\author{The Stacks Project Authors}
\date{stacks-project:homology.tex:local-snapshot-20261012; GNU FDL 1.2 or later}
\begin{document}
\maketitle
\noindent\textit{Editorial note.} All delivered mathematical text is the author's own, copied from the pinned Stacks Project chapter file homology.tex; no statement or proof step was rewritten, reordered or removed.\par\medskip
\noindent\textit{Editorial note.} The proof is the author's complete argument for all the implications, and each of them is written out inside the counted span. The statement opens with the author's own slogan environment, a comment-type environment of the chapter preamble, which is delivered as part of the counted statement.\par\medskip
\noindent\textit{Editorial note.} Every labelled reference occurring inside the counted spans is carried in assembly.external\_references as a non-load-bearing navigation pointer (same-chapter ones at their definition site, cross-chapter ones at the invocation inside the counted span), so no load-bearing mathematics is replaced by a stub. No bibliography citation occurs inside the counted spans.\par\medskip
\noindent\textit{Editorial note.} The same-name check for this card scanned every staged Stacks chapter of this repository -- all sources/stacks-*/pinned-original-source.tex files, the card's own chapter excluded -- and the labels found are recorded in dedup.same\_name\_cross\_chapter on the card.\par\medskip
\noindent\textit{Editorial note.} The preamble is the chapter's own native preamble excerpt already used for this source: the comment-type reference, slogan and history environments (native preamble lines 8-13), the Xy-pic setup (native preamble lines 14-18) and the native theorem declarations with the mathematical macros (native preamble lines 154-217, the actual last line of the file).\par\medskip
\noindent\textit{Allowed standard prerequisites (editorial).}\begin{itemize}
\item Exactness properties of adjoint functors. Categories, Lemma categories-lemma-exact-adjoint. Supplies that the left adjoint v is right exact, which together with the characterization of exactness gives the equivalence of (a) and (b), and is used again inside the two implications.
\item Characterizations of exact functors. Same chapter, Lemma lemma-exact-functor. Supplies the characterization of one-sided exactness through short exact sequences that turns injectivity-preservation of v into exactness of v.
\item Injective objects and the extension property. Same chapter, the definition of an injective object. Supplies the extension property used twice: once with I to extend a map out of vN, and once with u(I) to extend a map out of B.
\item Categories with enough injectives. Same chapter, the definition of a category with enough injectives. Supplies the monomorphism A -> I out of the kernel used in the last part of the proof.
\end{itemize}\medskip
\ifdefined\hypertarget\hypertarget{sourcepointer-1}{}\fi
\noindent\textit{Source reference.} Original lemma 0039: sources/stacks-categories/pinned-original-source.tex, source lines 3727--3727. Navigation pointer invoked at proof line 7506; the counted span contains the pointer only and the cited result is not reproduced inside it, so the pointer is registered as a non-load-bearing navigation stub in the established form for Stacks material. No bibliography citation is involved.\par\medskip
\ifdefined\hypertarget\hypertarget{sourcepointer-2}{}\fi
\noindent\textit{Source reference.} Original lemma 010N: sources/stacks-homology/pinned-original-source.tex, source lines 1762--1762. Navigation pointer invoked at proof line 7507; the counted span contains the pointer only and the cited result is not reproduced inside it, so the pointer is registered as a non-load-bearing navigation stub in the established form for Stacks material. No bibliography citation is involved.\par\medskip
\begin{lemma}
\label{lemma-adjoint-preserve-injectives}
\begin{slogan}
A functor with an exact left adjoint preserves injectives
\end{slogan}
Let $\mathcal{A}$ and $\mathcal{B}$ be abelian categories.
Let $u : \mathcal{A} \to \mathcal{B}$ and
$v : \mathcal{B} \to \mathcal{A}$ be additive functors with
$u$ right adjoint to $v$. Consider the following conditions:
\begin{enumerate}
\item[(a)] $v$ transforms injective maps into injective maps,
\item[(b)] $v$ is exact, and
\item[(c)] $u$ transforms injectives into injectives.
\end{enumerate}
Then (a) $\Leftrightarrow$ (b) $\Rightarrow$ (c). If $\mathcal{A}$
has enough injectives, then all three conditions are equivalent.
\end{lemma}

\begin{proof}
Observe that $v$ is right exact as a left adjoint (Categories, Lemma
\ref{categories-lemma-exact-adjoint}). Combined with
Lemma \ref{lemma-exact-functor} this explains why (a) $\Leftrightarrow$ (b).

\medskip\noindent
Assume (a). Let $I$ be an injective object of $\mathcal{A}$.
Let $\varphi : N \to M$ be an injective map in $\mathcal{B}$ and let
$\alpha : N \to uI$ be a morphism.
By adjointness we get a morphism $\alpha : vN \to I$ and
by assumption $v\varphi : vN \to vM$ is injective.
Hence as $I$ is an injective object we get a morphism
$\beta : vM \to I$ extending $\alpha$. By adjointness
again this corresponds to a morphism $\beta : M \to uI$ extending $\alpha$.
Hence (c) is true.

\medskip\noindent
Assume $\mathcal{A}$ has enough injectives and (c) holds.
Let  $f : B \to B'$ be a monomorphism in $\mathcal{B}$, and let
$A = \Ker(v(f))$. Choose a monomorphism $g : A \to I$ with $I$ injective
(it exists by assumption). Then $g$ extends to $g' : v(B) \to I$,
whence by adjunction a morphism $B \to u(I)$. Since $u(I)$ is injective,
this morphism extends to $h : B' \to u(I)$, hence by adjunction a morphism
$k : v(B') \to I$ extending $g'$. But then $k$ "extends'' $g$,
which forces $A = 0$ since $g$ was a monomorphism. Thus (a) is true.
\end{proof}

\end{document}
